\scanpage{27}
$p|_{\dot X_0\times I}$ induise $\dot X_0\times I\longrightarrow Y_0\times I$,
s'identifiant à $p_0\times\operatorname{id}_I$, où $p_0:X_0\longrightarrow Y_0$.

\deleted{Compatibilité}\inserted{Transversalité} d'une famille \deleted{finie}
$(B_i,E_{B_i})$ de voisinages cylindriques des parties $B_i$ de $X$.
La condition s'exprime en disant que pour
tout $J\subset I$, $\unsure{I|J}$ fini, la structure de
\[
\bigcap_{i\in J}E_{B_i}\overset{\mathrm{def}}=E_J,
\quad\text{autour de}\quad
B_J\overset{\mathrm{def}}=\bigcap_{i\in J}B_i,
\]
est celle de $B_J\times I^J$, la $\lambda$-homothétie
$\lambda_{E_j}$ de $E_j$ \inserted{(pour $j\in J$)} induisant
\[
(b,(t_\alpha)_{\alpha\in J})\longmapsto(b,(t'_\alpha)),
\qquad
t'_\alpha=t_\alpha\ \text{si }\alpha\ne j,\quad
\text{et }t'_\alpha=\lambda t_\alpha\ \text{si }\alpha=j.
\]
\margin{cet homéom. $E_J\simeq B_J\times I^J$ est
canonique, \unsure{via|si}
$\dot B_J=\bigcap_{j\in J}\dot E_j
\underset{\substack{\unsure{multiplié}\\\unsure{l'élément nul de}\\I^J}}{\simeq}B_J$
\unsure{\ldots}}

Cela peut aussi s'expliciter en exigeant
que $E_J$ soit cylindrique dans chaque $E_{B_j}$ etc\ldots
Cette relation de transversalité passe aux germes.

On dit que les familles précédentes \uline{et}
un voisinage cônique $T$ de $Y\subset X$ sont
\unsure{transversaux|transverses} \deleted{compatibles}
si \deleted{la structure des $T\cap E_i$}
\inserted{\deleted{de plus} $\forall\deleted{\unsure{j}}J$ comme dessus,}
\[
E_J\cap T\simeq\overline T_J\times I^J
\qquad(\text{où }\overline T_J=T\cap B_J)
\]
\margin{\unsure{s'induisant par} $E_J\simeq B_J\times I^J$ et}
induisant
\[
E_J\cap\dot T\simeq\dot T_J\times I^J
\qquad(\text{où }\dot T_J=\dot T\cap B_J)
\]
\[
E_J\cap Y\simeq Y_J\times I^J
\qquad(\text{où }Y_J=Y\cap B_J)
\]
\margin{26}
